\subsection{Closing the loop among meal plans, recipes, and groceries}
The proposal promises a grocery list grounded in planned meals. Without aggregation from \texttt{MealPlan} and \texttt{Recipe} ingredients, users still bridge planning and shopping mentally. We are defining rules for servings scaling, deduplication, and unit handling before shipping a ``sync from plan'' action, and will preserve full edit control after generation.

\subsection{Search, filter, and dietary personalization}
Filtering by dietary restrictions, cook time, and difficulty requires both data (tags or enums on recipes) and UI tied to \texttt{users.me} preferences. We plan debounced search and high-signal filters (e.g., dietary tags) to limit choice overload on Home~\cite{hartson2012ux}.

\subsection{Mise en place and step granularity}
Showing all ingredients on every step can overwhelm; showing none undermines guided cooking. We will extend the schema to associate ingredients with steps (or infer links from parsing), surface a short checklist for the active step only, and compare error rates and subjective load against detail-only ingredient views in testing.

\subsection{Trustworthy AI parsing}
A mock parser can mislead participants about accuracy. We will swap in a real text parser behind the same verify-and-edit screen, keep image/video import clearly labeled as beta or upcoming, and measure correction time and satisfaction~\cite{nielsen1994heuristic}.

\subsection{Calendar information density}
A month grid supports an ``at a glance'' metaphor but clutters with multiple meals per day. We will prototype indicators on a month view layered on the current day-detail pattern and test week vs.\ default month and the number of taps to add a recipe from the library.

\subsection{Engineering risks (brief)}
Recipe list endpoints currently return a global set; grocery item updates must be scoped by list ownership in future hardening. These items are tracked in the project backlog (\texttt{SimpleChef/CONTINUATION\_CHECKLIST.md}).
